\hspace{3mm}

\centerline{\bf \Huge  Задачи по графам.}
\centerline{\bf \Huge Уровень 2: Древний Рус}

\begin{flushright}
    \scriptsize{Выпьем воды Байкальской братья!\\ За победу над ящером окоянным.}
\end{flushright}

Связный граф без циклов называется \textit{деревом.}

\problem{0} Несколько слов про деревья

\noindent a) Пусть в графе нет циклов и вершин степени ноль(то есть \textit{изолированная} вершина). Докажите, что есть вершина степени один. Такая вершина, кстати, ещё называется \textit{висячей.}

 \noindent b) Докажите, что их на самом деле хотя бы две.

 \noindent c) Пусть в дереве есть вершина степени d. Докажите, что висячих вершин на самом деле хотя бы d.

\textbf{Великорусское подвешивание графа за вершину.} Назовём подвешиванием (ранжированием) графа за вершину следующую процедуру. Выберем любую вершину графа. Назовём её нулевым уровнем (рангом). На первом уровне расположим все вершины, соединённые с вершиной нулевого уровня. На каждом следующем уровне будем располагать все вершины, соединённые с вершинами предыдущего уровня и не расположенные на более ранних уровнях. Процесс продолжается, пока все вершины не расположатся на уровнях.

\problem{1} Докажите, что в любом дереве есть вершина степени 1.

\problem{2} Докажите, что в любом связном графе найдётся вершина, при удалении которой(вместе с исходящими из нее ребрами) оставшийся граф будет связным.

\problem{3} На открытии ЛМШ 2024 каждый ученик пожал руку не более чем четырём другим ученикам. Также известно, что любой ученик может познакомиться с любым другим не более чем через 2 новых рукопожатия (например, если Герман пожал руку Баиру, Баир Глебу, а Глеб Дольгану, а Дольган, то Германа и Ердема разделяет два рукопожатия $-$ с Глебом и Долей). Докажите, что на открытии ЛМШ 2024 было не более 53 ученика.

\problem{4} Степень всех вершин графа не меньше $n$, причем в нем нет циклов длины $3, 4, 5$. Докажите, что в нем найдутся $n^2 - n$ вершин, никакие две из которых не смежны.

\problem{5} В стране 100 городов, некоторые из которых соединены авиалиниями. Известно, что от каждого города можно долететь до любого другого (возможно, с пересадками). Докажите, что можно побывать во всех городах, совершив не более 196 перелётов.